\documentclass[11pt]{article}
\usepackage[margin=1in]{geometry}
\usepackage{amsmath,amssymb,amsthm,enumitem,hyperref}
\hypersetup{colorlinks=true,linkcolor=blue,urlcolor=blue,citecolor=blue}
\newtheorem{theorem}{Theorem}
\newtheorem{lemma}{Lemma}
\newcommand{\E}{\mathbb E}
\newcommand{\Q}{\mathcal Q}
\newcommand{\Z}{\mathbb Z}
\title{Bivariate Ehrhart reciprocity for preorder polytopes\\\large MathDB 374965}
\author{Principia Math In Progress}
\date{September 18, 2026}
\begin{document}
\maketitle
\begin{abstract}
We prove Conjecture 7.2 of Athanasiadis and Chapoton: the double
Ehrhart polynomial of every finite preorder satisfies
$\E_\tau(-u,-v)=(-1)^n\E_\tau(u-1,v-1)$.
The proof translates interior lattice points by the all-ones vector
and applies ordinary Ehrhart reciprocity on each positive integral ray.
\end{abstract}

\section*{Task 1. Full solution}
Let $\tau$ be a preorder on a nonempty finite set $E$, with $|E|=n$.
An order ideal $I$ satisfies $b\in I$, $a\le_\tau b\Rightarrow a\in I$.
This definition permits equivalent distinct elements of the preorder.
For $I\subseteq E$, write
$\Delta_I=\operatorname{conv}(0,e_i:i\in I)\subseteq\mathbb R^E$,
and let $U(e)$ denote the principal order filter generated by $e$.
The two-parameter polytope of \cite[Section 3.2]{AC} is
\[
 \Q_\tau(r,s)=r\sum_{e\in E}\Delta_{U(e)}+s\Delta_E,
 \qquad r,s\in\Z_{\ge0}.
\]
Summing once for each element is the same as weighting the summand
for an equivalence class by its cardinality, as in the source.
By \cite[Proposition 3.3]{AC}, its inequality description is
\begin{equation}\label{eq:ineq}
 x_e\ge0\quad(e\in E),\qquad
 \sum_{e\in I}x_e\le r|I|+s
 \quad(\varnothing\ne I\text{ an order ideal}).
\end{equation}
It is a lattice polytope, being an integral Minkowski sum of lattice
simplices. The polynomial $\E_\tau(u,v)$ of total degree $n$ is
characterized by
\[
 \E_\tau(r,s)=\#(\Q_\tau(r,s)\cap\Z^E)
 \quad(r,s\in\Z_{\ge0});
\]
its existence is established in \cite[Proposition 7.1]{AC}.

\begin{theorem}
For every finite preorder $\tau$ of size $n$, the following is a
polynomial identity:
\begin{equation}\label{eq:target}
 \boxed{\E_\tau(-u,-v)=(-1)^n\E_\tau(u-1,v-1).}
\end{equation}
\end{theorem}

\begin{lemma}\label{lem:shift}
For positive integers $r,s$, translation by $-\mathbf1$, where
$\mathbf1=(1)_{e\in E}$, is a bijection
\[
 \operatorname{int}\Q_\tau(r,s)\cap\Z^E
 \longrightarrow \Q_\tau(r-1,s-1)\cap\Z^E.
\]
\end{lemma}
\begin{proof}
The polytope $\Q_\tau(r,s)$ is full dimensional: it contains
$[0,r]^E$ by \eqref{eq:ineq}. Every normal in \eqref{eq:ineq} is
nonzero. Thus its interior consists exactly of the points satisfying
all those inequalities strictly. This remains true for redundant
inequalities: equality with a valid nonzero normal cannot occur at
an interior point, and strict satisfaction of a finite system gives
an open neighborhood contained in the polytope.

For an integral point $x$, the strict inequalities say precisely
\[
 x_e\ge1,
 \qquad \sum_{e\in I}x_e\le r|I|+s-1
 \quad(\varnothing\ne I\text{ an order ideal}).
\]
Set $y=x-\mathbf1$. The same conditions become
\[
 y_e\ge0,
 \qquad \sum_{e\in I}y_e\le(r-1)|I|+(s-1).
\]
By \eqref{eq:ineq}, these are exactly the defining inequalities for
$\Q_\tau(r-1,s-1)$. Both translations preserve the integer lattice,
which proves the bijection, including $r=1$ or $s=1$.
\end{proof}

\begin{proof}[Proof of the theorem]
Fix positive integers $r,s$, and put $P=\Q_\tau(r,s)$.
For every nonnegative integer $k$ the Minkowski definition gives
$kP=\Q_\tau(kr,ks)$. Hence the ordinary Ehrhart polynomial
$L_P(k)=\#(kP\cap\Z^E)$ satisfies
\[
 L_P(k)=\E_\tau(kr,ks)
\]
as an identity of polynomials in $k$, since the two polynomials
agree at all nonnegative integers. In particular this equality is
valid at $k=-1$. Ordinary Ehrhart--Macdonald reciprocity, applicable
to the full dimensional lattice polytope $P$, now gives
\[
 \E_\tau(-r,-s)=L_P(-1)
 =(-1)^n\#(\operatorname{int}P\cap\Z^E)
 =(-1)^n\E_\tau(r-1,s-1),
\]
where the last equality is Lemma~\ref{lem:shift}.

The difference between the two sides of \eqref{eq:target} is a
polynomial in $u,v$ vanishing at every positive integer pair.
For each fixed positive integer $v$, it has infinitely many zeros
as a polynomial in $u$, so all its coefficients in $u$ vanish at
that $v$. Each such coefficient is a polynomial in $v$ vanishing
at all positive integers and is therefore identically zero.
This proves \eqref{eq:target} for arbitrary $u,v$.
\end{proof}

If the empty preorder is allowed, its polynomial is $1$, and the
identity holds for $n=0$ as well. Values with zero or negative
parameters in the theorem are polynomial evaluations; the proof
does not require giving negatively scaled Minkowski sums a counting
interpretation.

\subsection*{Attribution, scope, and verification}
The polytope descriptions and the existence of the double Ehrhart
polynomial are established results of Athanasiadis and Chapoton.
The contribution here is the translation argument and its application
to their Conjecture 7.2. Their Remark 7.3 already explains how this
identity implies their conjectured descent formula.
Ordinary Ehrhart reciprocity is a classical theorem; no multivariate
reciprocity theorem is needed.

The exact verifier independently enumerates Minkowski sums and
inequality-defined lattice-point sets for all labeled preorders on
at most four elements. It also checks the interior translation, the
full polynomial identity using the source's draconian formula, and
the published five-element example. These finite computations
supplement the proof for arbitrary size.

All 26 pages of the current primary were read. Bounded public
checks found no earlier full claim for this target and no identified
journal final. The related paper \cite{D} proves support-enumerator
duality, a different statement; its relevant scope was checked.
Review is by the authoring agent only, without independent peer
review or proof-assistant certification.

\section*{Tasks 2--5. Proposed follow-up directions}
\begin{enumerate}[label=\textbf{Task \arabic*.},start=2,leftmargin=*]
\item Give a direct word bijection for the descent formula implied
by this result and \cite[Remark 7.3]{AC}, avoiding polynomial
reciprocity in that combinatorial interpretation.
\item Classify multivariate Minkowski families with a fixed integral
translation carrying all interior lattice points to a family with
the parameters decreased by prescribed integers.
\item Refine the translation by coordinate monomial weights and
derive a reciprocity formula for multivariate generating functions.
\item Determine the restrictions on the mixed coefficients of
$\E_\tau$ obtained from the parity of
$\E_\tau(u-\tfrac12,v-\tfrac12)$ under $(u,v)\mapsto(-u,-v)$.
\end{enumerate}
\begin{thebibliography}{9}
\raggedright
\bibitem{AC} C. A. Athanasiadis and F. Chapoton,
\emph{Polytopes and posets associated to preorders},
\href{https://arxiv.org/abs/2605.26916v1}{arXiv:2605.26916v1}, 2026.
\bibitem{D} Z. Dai, Q. Hou, Z. Liu, W. Thawinrak, and H. Wang,
\emph{Counting Lattice Points in Minkowski Sums of Cross Polytopes},
\href{https://arxiv.org/abs/2608.16037v2}{arXiv:2608.16037v2}, 2026.
\end{thebibliography}
\end{document}
